\subsection{Проблема восстановления данных и классическая интерполяция}Применение дискретного вейвлет-преобразования с последующим отсечением высокочастотных поддиапазонов позволяет достичь значительной степени сжатия, однако полученный низкочастотный аппроксимирующий сигнал (LL-слой) имеет геометрическое разрешение, уменьшенное в два раза по каждой из координатных осей \cite{mallat2005}. Для демонстрации или дальнейшей обработки такого изображения на стороне приемника возникает математически некорректная обратная задача — восстановление исходного пространственного разрешения (апскейлинг) на основе ограниченного набора данных.Традиционные подходы к решению данной задачи базируются на алгоритмах пространственной интерполяции, таких как билинейная или бикубическая интерполяция \cite{gonzalez2012}. Бикубическая интерполяция вычисляет значение нового пикселя как взвешенную сумму 16 ближайших соседних пикселей в оригинальной сетке, используя кубические сплайны для обеспечения гладкости функции интенсивности.Несмотря на вычислительную простоту и повсеместное использование, классические интерполяционные фильтры обладают фундаментальным ограничением: они функционируют как фильтры нижних частот (low-pass filters), сглаживая перепады яркости \cite{bovik2009}. Интерполяция физически не способна восстановить высокочастотную информацию (мелкие текстуры, резкие грани объектов), которая была утрачена на этапе сжатия. При применении бикубической интерполяции к сильно сжатому вейвлет-скелету результирующее изображение получается чрезмерно размытым, лишенным микроконтраста и деталей, что делает данный метод неприемлемым для качественной реконструкции графических данных \cite{richardson2005}.\subsection{Технологии нейросетевого сверхразрешения}Преодоление фундаментальных ограничений классической интерполяции стало возможным с развитием методов глубокого машинного обучения, в частности, технологий однокадрового сверхвысокого разрешения (Single Image Super-Resolution, SISR). Задача SISR формулируется как поиск функции нелинейного отображения, способной предсказывать недостающие высокочастотные детали и реконструировать изображение высокого разрешения (High Resolution, HR) из его версии низкого разрешения (Low Resolution, LR) \cite{dong2015srcnn}.Сверточные нейронные сети (Convolutional Neural Networks, CNN) позволили перевести задачу апскейлинга из плоскости чисто математического усреднения в плоскость интеллектуальной генерации признаков. Пионерская работа в этой области — архитектура SRCNN (Super-Resolution Convolutional Neural Network) — продемонстрировала, что глубокая сеть способна обучаться сложным (end-to-end) преобразованиям между LR и HR патчами изображений, значительно превосходя классические фильтры \cite{dong2015srcnn}.Для объективной оценки качества восстановления в технологиях сверхразрешения традиционно применяются метрики, моделирующие как математическую точность, так и человеческое восприятие \cite{salomon2004}:\begin{itemize}\item \textbf{PSNR (Пиковое отношение сигнала к шуму)} — оценивает абсолютную среднеквадратичную ошибку между оригиналом и восстановленным изображением. Измеряется в децибелах (дБ).\item \textbf{SSIM (Индекс структурного сходства)} — оценивает деградацию структурной информации, изменения контраста и яркости. Значение метрики варьируется в диапазоне $[-1, 1]$, где $1$ означает абсолютную идентичность.\end{itemize}\subsection{Архитектура нейронной сети EDSR}Современным стандартом в задачах SISR, демонстрирующим оптимальный баланс между качеством восстановления и вычислительной стабильностью, является архитектура EDSR (Enhanced Deep Residual Networks), предложенная B. Lim и др. \cite{lim2017edsr}. Данная модель стала победителем престижного соревнования NTIRE 2017, доказав свое превосходство над предшествующими архитектурами, такими как SRResNet. Именно модель EDSR с коэффициентом масштабирования $\times 2$ интегрирована в разработанный гибридный программный комплекс для восстановления изображений после вейвлет-сжатия.Ключевым архитектурным нововведением EDSR, обусловившим её высокую эффективность, стал радикальный пересмотр структуры остаточных блоков (Residual Blocks). В классических глубоких сетях для классификации изображений (например, ResNet) после каждого сверточного слоя применяется слой пакетной нормализации (Batch Normalization, BN). Авторы EDSR доказали, что в задачах сверхразрешения пакетная нормализация наносит вред: она ограничивает диапазон значений признаков (избавляясь от абсолютных значений яркости и контраста) и потребляет значительный объем оперативной памяти \cite{lim2017edsr}.Удаление слоев BN из остаточных блоков позволило:\begin{enumerate}\item Сохранить абсолютные значения интенсивностей пикселей, что критически важно для генерации реалистичных цветов и текстур.\item Существенно снизить потребление памяти, что дало возможность увеличить глубину сети (количество слоев) и ширину (количество фильтров в свертках), повысив общую выразительную способность модели без дополнительных вычислительных затрат \cite{lim2017edsr}.\end{enumerate}В контексте разрабатываемой системы гибридного сжатия, применение сети EDSR решает ключевую задачу: модель, принимая на вход нормализованный низкочастотный вейвлет-скелет, не просто масштабирует его, а применяет усвоенные в процессе предварительного обучения статистические закономерности для генерации правдоподобных высокочастотных компонент, фактически восстанавливая информацию, отброшенную на этапе DWT-декомпозиции.